%% ============================================================================
%% CAPÍTULO 6: EL ALGORITMO DE SIMPLIFICACIÓN
%% ============================================================================
\chapter{Contribución III: Un enfoque alternativo para simplificación de grafos que preserva topología}

\section{Diseño del Motor de Simplificación ACJ}
% GUÍA: Visión general del algoritmo propuesto que evita la destrucción geométrica y la cristalización.

\section{Fase 1: Filtro Topológico-Vectorial}
% GUÍA: Explicar el uso de álgebra lineal (producto punto y cálculo de ángulos de desviación) para empaquetar calles de grado 2 sin borrar curvas críticas o "esquinas peligrosas".

\section{Fase 2: Simplificación Dinámica basada en Voronoi}
% GUÍA: Explicar la aplicación de Douglas-Peucker con una tolerancia (épsilon) dinámica calculada a partir del área local de teselación de Voronoi.
% Mencionar la prevención de colisiones usando índices espaciales (STRtree / R-Trees).

\section{Formulación de Métricas de Evaluación Propias}
% GUÍA: Presentar las métricas matemáticas creadas para medir el éxito de la simplificación:
% 1. Ratio de Compresión Topológica ($|V_{simp}| / |V_{raw}|$).
% 2. Error de Invarianza de Longitud de Aristas ($\Delta L / L_{raw}$).
% 3. Grado de preservación de conectividad.